\documentclass[aspectratio=1610,11pt]{beamer}
% Shared style of the lecture slides (Slides/slidesNN.tex). Each deck starts with
%     \documentclass[aspectratio=1610,11pt]{beamer}
%     % Shared style of the lecture slides (Slides/slidesNN.tex). Each deck starts with
%     \documentclass[aspectratio=1610,11pt]{beamer}
%     % Shared style of the lecture slides (Slides/slidesNN.tex). Each deck starts with
%     \documentclass[aspectratio=1610,11pt]{beamer}
%     \input{slides_preamble}
% aspectratio=1610 (16:10) matches the Samsung Galaxy Tab S9 FE+ (2560 x 1600) used for teaching.
%
% The decks have no overlays (\pause, \only): every slide is one PDF page, so what you write on a
% page while sharing the tablet screen stays on that page. Board pages (\board) are white pages with
% a small header, for calculations and proofs written live.

\usepackage[T1]{fontenc}
\usepackage{lmodern}
\usefonttheme[onlymath]{serif}
\usepackage{amsmath,amssymb}
\usepackage{booktabs}
\usepackage{tikz}
\usetikzlibrary{arrows.meta,calc,positioning,shapes.geometric,decorations.pathreplacing}
\usepackage{pgfplots}
\pgfplotsset{compat=1.18}
\usepgfplotslibrary{fillbetween}

% Numbers of the lecture notes. A deck says which notes it follows with \setnotes{2} and
% loads their labels with \externaldocument[L2-]{../Lecture2/lecture2} (the .aux of the notes
% must exist). Then \begin{theorem}[Spectral theorem][thm:spectral] prints "notes L2, 10.1" at
% the right of the block title, and \secref{sec:spectral} prints "§2.10".
\usepackage{xr-hyper}
\makeatletter
\newcommand{\setnotes}[1]{\def\slides@notes{#1}}
\setnotes{2}
\newcommand{\notesref}[1]{\ref*{L\slides@notes-#1}}
\newcommand{\secref}[1]{\S\slides@notes.\ref*{L\slides@notes-#1}}
\makeatother

% Macros of the lecture notes.
\newcommand{\norm}[1]{\lVert #1 \rVert}
\newcommand{\ip}[2]{\langle #1, #2 \rangle}
\newcommand{\py}[1]{\texttt{\detokenize{#1}}}
\newcommand{\E}{\mathbb{E}}
\newcommand{\indep}{\perp\!\!\!\perp}
\newcommand{\R}{\mathbb{R}}
\newcommand{\bx}{\mathbf{x}}
\newcommand{\by}{\mathbf{y}}
\newcommand{\bw}{\mathbf{w}}
\newcommand{\bv}{\mathbf{v}}
\newcommand{\bu}{\mathbf{u}}
\newcommand{\bq}{\mathbf{q}}
\newcommand{\bb}{\mathbf{b}}
\newcommand{\be}{\mathbf{e}}
\DeclareMathOperator{\Var}{Var}
\DeclareMathOperator{\Cov}{Cov}
\DeclareMathOperator{\Corr}{Corr}
\DeclareMathOperator{\KL}{KL}
\DeclareMathOperator{\im}{im}
\DeclareMathOperator{\rank}{rank}
\DeclareMathOperator{\Null}{null}
\DeclareMathOperator{\col}{col}
\DeclareMathOperator{\row}{row}
\DeclareMathOperator{\Span}{span}
\DeclareMathOperator{\tr}{tr}
\DeclareMathOperator{\diag}{diag}
\DeclareMathOperator{\Div}{div}
\DeclareMathOperator{\curl}{curl}
\DeclareMathOperator{\vol}{vol}
\DeclareMathOperator{\area}{area}
\DeclareMathOperator{\sign}{sign}
\DeclareMathOperator{\softmax}{softmax}
\DeclareMathOperator{\ReLU}{ReLU}
\DeclareMathOperator*{\argmin}{arg\,min}
\DeclareMathOperator*{\argmax}{arg\,max}

% Colors.
\definecolor{accent}{RGB}{31,78,140}      % titles, theorems
\definecolor{defcolor}{RGB}{24,110,92}    % definitions
\definecolor{excolor}{RGB}{150,82,20}     % examples
\definecolor{mlcolor}{RGB}{120,40,120}    % machine learning connections
\definecolor{boardink}{RGB}{140,140,140}  % header of the board pages

% Layout.
\setbeamertemplate{navigation symbols}{}
\setbeamersize{text margin left=0.9em, text margin right=0.9em}
\setbeamercolor{structure}{fg=accent}
\setbeamercolor{normal text}{fg=black, bg=white}
\setbeamercolor{frametitle}{fg=accent, bg=white}
\setbeamerfont{frametitle}{series=\bfseries, size=\large}
\setbeamertemplate{frametitle}{%
    \vskip0.45em\insertframetitle\par\vskip-0.25em
    {\color{accent!25}\rule{\linewidth}{0.5pt}}\vskip-0.2em}
\setbeamertemplate{itemize item}{\color{accent}\raisebox{0.15ex}{\scriptsize$\blacktriangleright$}}
\setbeamertemplate{itemize subitem}{\color{accent}\raisebox{0.2ex}{\tiny$\bullet$}}
\setbeamertemplate{enumerate item}{\color{accent}\insertenumlabel.}
\setbeamerfont{footline}{size=\tiny}
\setbeamertemplate{footline}{%
    \leavevmode\hbox to \paperwidth{\hskip0.9em{\color{gray}\insertshorttitle}\hfill
    {\color{gray}\insertframenumber}\hskip0.9em}\vskip0.45em}
\setbeamertemplate{blocks}[rounded][shadow=false]
\setbeamercolor{block title}{fg=white, bg=accent}
\setbeamercolor{block body}{fg=black, bg=accent!6}
\setbeamerfont{block title}{size=\normalsize, series=\bfseries}
\setlength{\abovedisplayskip}{4pt}
\setlength{\belowdisplayskip}{4pt}

% Blocks: \begin{theorem}[Spectral theorem][thm:spectral] ... \end{theorem}; both arguments are optional.
\makeatletter
\newcommand{\slides@block}[5]{% color, word, name, label in the notes, body tint
    \setbeamercolor{block title}{fg=white, bg=#1}%
    \setbeamercolor{block body}{fg=black, bg=#1!#5}%
    \begin{block}{#2\if\relax\detokenize{#3}\relax\else\ {\normalfont(#3)}\fi
        \if\relax\detokenize{#4}\relax\else\hfill{\normalfont\scriptsize notes~L\slides@notes, \notesref{#4}}\fi}}
\DeclareDocumentEnvironment{theorem}{O{} O{}}{\slides@block{accent}{Theorem}{#1}{#2}{6}}{\end{block}}
\DeclareDocumentEnvironment{proposition}{O{} O{}}{\slides@block{accent}{Proposition}{#1}{#2}{6}}{\end{block}}
\DeclareDocumentEnvironment{lemma}{O{} O{}}{\slides@block{accent}{Lemma}{#1}{#2}{6}}{\end{block}}
\DeclareDocumentEnvironment{corollary}{O{} O{}}{\slides@block{accent}{Corollary}{#1}{#2}{6}}{\end{block}}
\DeclareDocumentEnvironment{definition}{O{} O{}}{\slides@block{defcolor}{Definition}{#1}{#2}{6}}{\end{block}}
\DeclareDocumentEnvironment{example}{O{} O{}}{\slides@block{excolor}{Example}{#1}{#2}{6}}{\end{block}}
\DeclareDocumentEnvironment{remark}{O{} O{}}{\slides@block{gray!70!black}{Remark}{#1}{#2}{5}}{\end{block}}
\DeclareDocumentEnvironment{mlbox}{O{In machine learning}}{%
    \setbeamercolor{block title}{fg=white, bg=mlcolor}%
    \setbeamercolor{block body}{fg=black, bg=mlcolor!6}%
    \begin{block}{#1}}{\end{block}}
\makeatother

% Board pages.
%   \board{Title}{Task, shown in small print at the top}            one white page
%   \board[3]{Title}{Task}                                          three pages (the last two blank)
\newcounter{boardpage}
\newcommand{\boardhead}[1]{%
    \vskip0.3em
    \noindent{\color{boardink}\scriptsize
        \tikz[baseline=(b.base)]\node[draw=boardink, rounded corners=2pt, inner sep=1.5pt,
            font=\tiny\bfseries\sffamily] (b) {BOARD};\ \ #1}\par\vskip0.3em}
\newcommand{\board}[3][1]{%
    \begin{frame}[t]
        \boardhead{#2}
        \if\relax\detokenize{#3}\relax\else
            {\small #3\par}
        \fi
    \end{frame}%
    \setcounter{boardpage}{1}%
    \loop\ifnum\value{boardpage}<#1\relax
        \stepcounter{boardpage}%
        \begin{frame}[t]\boardhead{#2 (continued)}\end{frame}%
    \repeat}

% Practice pages: problems of the practice file, as white pages for solving them live.
% A deck loads the practice file's labels with \externaldocument[P1-]{../Lecture1/practice1} and
% \setpractice{1}. Then
%   \practice{pr:eig2}{\mandatory}{Problem text}    one page, headed "PRACTICE  Problem 10.1 (mandatory)"
%   \practice[2]{pr:eig-props}{}{Problem text}      two pages (the second blank)
%   \practice{pr:eig-gre}{$\star$}{Problem text}    the marks of the practice file
% The answers stay in LectureN/practiceN_answers.tex.
\newcommand{\mandatory}{\textcolor{red!70!black}{\textbf{mandatory}}}
\makeatletter
\newcommand{\setpractice}[1]{\def\slides@practice{#1}}
\newcommand{\practiceref}[1]{\ref*{P\slides@practice-#1}}
\newcommand{\practicehead}[2]{% label in the practice file, text after the problem number
    \vskip0.3em
    \noindent{\color{boardink}\scriptsize
        \tikz[baseline=(b.base)]\node[draw=defcolor, text=defcolor, rounded corners=2pt, inner sep=1.5pt,
            font=\tiny\bfseries\sffamily] (b) {PRACTICE};\ \ Practice~\slides@practice, Problem~\practiceref{#1}#2}%
    \par\vskip0.3em}
\newcommand{\practice}[4][1]{%
    \begin{frame}[t]
        \practicehead{#2}{\if\relax\detokenize{#3}\relax\else\ (#3)\fi}
        {\small #4\par}
    \end{frame}%
    \setcounter{boardpage}{1}%
    \loop\ifnum\value{boardpage}<#1\relax
        \stepcounter{boardpage}%
        \begin{frame}[t]\practicehead{#2}{\ (continued)}\end{frame}%
    \repeat}
\makeatother

% Divider page: \topic{\secref{sec:spectral}}{The spectral theorem} or \topic{\S2.1--2.2}{Norms}.
\newcommand{\topic}[2]{%
    \begin{frame}[c]
        \centering
        {\color{gray}\large #1}\par\vskip0.4em
        {\color{accent}\LARGE\bfseries #2}\par
    \end{frame}}

% Title page.
\setbeamercolor{title}{fg=accent}
\setbeamerfont{title}{series=\bfseries, size=\Large}
\setbeamercolor{subtitle}{fg=gray!60!black}
\setbeamerfont{subtitle}{size=\normalsize}
\setbeamertemplate{title page}{%
    \vfill
    {\usebeamercolor[fg]{title}\usebeamerfont{title}\inserttitle\par}
    \vskip0.6em
    {\usebeamercolor[fg]{subtitle}\usebeamerfont{subtitle}\insertsubtitle\par}
    \vskip1.2em
    {\small\color{gray}\insertauthor\par}
    \vfill}
\author{Mathematics for Machine Learning}
\date{}

% Two columns with a picture: \sidebyside{0.5}{left text}{right picture}.
\newcommand{\sidebyside}[3]{%
    \begin{columns}[T]
        \begin{column}{#1\linewidth}#2\end{column}
        \begin{column}{\dimexpr\linewidth-#1\linewidth-1em\relax}#3\end{column}
    \end{columns}}

% "Check" slide after a board computation.
\newcommand{\checkslide}[2]{%
    \begin{frame}{Check: #1}
        #2
    \end{frame}}

% Pictures.
\newcommand{\notesfig}[2][1]{\resizebox{#1\linewidth}{!}{\input{figures/notes/#2}}}
\newcommand{\fig}[2][1]{\includegraphics[width=#1\linewidth]{figures/#2}}

% aspectratio=1610 (16:10) matches the Samsung Galaxy Tab S9 FE+ (2560 x 1600) used for teaching.
%
% The decks have no overlays (\pause, \only): every slide is one PDF page, so what you write on a
% page while sharing the tablet screen stays on that page. Board pages (\board) are white pages with
% a small header, for calculations and proofs written live.

\usepackage[T1]{fontenc}
\usepackage{lmodern}
\usefonttheme[onlymath]{serif}
\usepackage{amsmath,amssymb}
\usepackage{booktabs}
\usepackage{tikz}
\usetikzlibrary{arrows.meta,calc,positioning,shapes.geometric,decorations.pathreplacing}
\usepackage{pgfplots}
\pgfplotsset{compat=1.18}
\usepgfplotslibrary{fillbetween}

% Numbers of the lecture notes. A deck says which notes it follows with \setnotes{2} and
% loads their labels with \externaldocument[L2-]{../Lecture2/lecture2} (the .aux of the notes
% must exist). Then \begin{theorem}[Spectral theorem][thm:spectral] prints "notes L2, 10.1" at
% the right of the block title, and \secref{sec:spectral} prints "§2.10".
\usepackage{xr-hyper}
\makeatletter
\newcommand{\setnotes}[1]{\def\slides@notes{#1}}
\setnotes{2}
\newcommand{\notesref}[1]{\ref*{L\slides@notes-#1}}
\newcommand{\secref}[1]{\S\slides@notes.\ref*{L\slides@notes-#1}}
\makeatother

% Macros of the lecture notes.
\newcommand{\norm}[1]{\lVert #1 \rVert}
\newcommand{\ip}[2]{\langle #1, #2 \rangle}
\newcommand{\py}[1]{\texttt{\detokenize{#1}}}
\newcommand{\E}{\mathbb{E}}
\newcommand{\indep}{\perp\!\!\!\perp}
\newcommand{\R}{\mathbb{R}}
\newcommand{\bx}{\mathbf{x}}
\newcommand{\by}{\mathbf{y}}
\newcommand{\bw}{\mathbf{w}}
\newcommand{\bv}{\mathbf{v}}
\newcommand{\bu}{\mathbf{u}}
\newcommand{\bq}{\mathbf{q}}
\newcommand{\bb}{\mathbf{b}}
\newcommand{\be}{\mathbf{e}}
\DeclareMathOperator{\Var}{Var}
\DeclareMathOperator{\Cov}{Cov}
\DeclareMathOperator{\Corr}{Corr}
\DeclareMathOperator{\KL}{KL}
\DeclareMathOperator{\im}{im}
\DeclareMathOperator{\rank}{rank}
\DeclareMathOperator{\Null}{null}
\DeclareMathOperator{\col}{col}
\DeclareMathOperator{\row}{row}
\DeclareMathOperator{\Span}{span}
\DeclareMathOperator{\tr}{tr}
\DeclareMathOperator{\diag}{diag}
\DeclareMathOperator{\Div}{div}
\DeclareMathOperator{\curl}{curl}
\DeclareMathOperator{\vol}{vol}
\DeclareMathOperator{\area}{area}
\DeclareMathOperator{\sign}{sign}
\DeclareMathOperator{\softmax}{softmax}
\DeclareMathOperator{\ReLU}{ReLU}
\DeclareMathOperator*{\argmin}{arg\,min}
\DeclareMathOperator*{\argmax}{arg\,max}

% Colors.
\definecolor{accent}{RGB}{31,78,140}      % titles, theorems
\definecolor{defcolor}{RGB}{24,110,92}    % definitions
\definecolor{excolor}{RGB}{150,82,20}     % examples
\definecolor{mlcolor}{RGB}{120,40,120}    % machine learning connections
\definecolor{boardink}{RGB}{140,140,140}  % header of the board pages

% Layout.
\setbeamertemplate{navigation symbols}{}
\setbeamersize{text margin left=0.9em, text margin right=0.9em}
\setbeamercolor{structure}{fg=accent}
\setbeamercolor{normal text}{fg=black, bg=white}
\setbeamercolor{frametitle}{fg=accent, bg=white}
\setbeamerfont{frametitle}{series=\bfseries, size=\large}
\setbeamertemplate{frametitle}{%
    \vskip0.45em\insertframetitle\par\vskip-0.25em
    {\color{accent!25}\rule{\linewidth}{0.5pt}}\vskip-0.2em}
\setbeamertemplate{itemize item}{\color{accent}\raisebox{0.15ex}{\scriptsize$\blacktriangleright$}}
\setbeamertemplate{itemize subitem}{\color{accent}\raisebox{0.2ex}{\tiny$\bullet$}}
\setbeamertemplate{enumerate item}{\color{accent}\insertenumlabel.}
\setbeamerfont{footline}{size=\tiny}
\setbeamertemplate{footline}{%
    \leavevmode\hbox to \paperwidth{\hskip0.9em{\color{gray}\insertshorttitle}\hfill
    {\color{gray}\insertframenumber}\hskip0.9em}\vskip0.45em}
\setbeamertemplate{blocks}[rounded][shadow=false]
\setbeamercolor{block title}{fg=white, bg=accent}
\setbeamercolor{block body}{fg=black, bg=accent!6}
\setbeamerfont{block title}{size=\normalsize, series=\bfseries}
\setlength{\abovedisplayskip}{4pt}
\setlength{\belowdisplayskip}{4pt}

% Blocks: \begin{theorem}[Spectral theorem][thm:spectral] ... \end{theorem}; both arguments are optional.
\makeatletter
\newcommand{\slides@block}[5]{% color, word, name, label in the notes, body tint
    \setbeamercolor{block title}{fg=white, bg=#1}%
    \setbeamercolor{block body}{fg=black, bg=#1!#5}%
    \begin{block}{#2\if\relax\detokenize{#3}\relax\else\ {\normalfont(#3)}\fi
        \if\relax\detokenize{#4}\relax\else\hfill{\normalfont\scriptsize notes~L\slides@notes, \notesref{#4}}\fi}}
\DeclareDocumentEnvironment{theorem}{O{} O{}}{\slides@block{accent}{Theorem}{#1}{#2}{6}}{\end{block}}
\DeclareDocumentEnvironment{proposition}{O{} O{}}{\slides@block{accent}{Proposition}{#1}{#2}{6}}{\end{block}}
\DeclareDocumentEnvironment{lemma}{O{} O{}}{\slides@block{accent}{Lemma}{#1}{#2}{6}}{\end{block}}
\DeclareDocumentEnvironment{corollary}{O{} O{}}{\slides@block{accent}{Corollary}{#1}{#2}{6}}{\end{block}}
\DeclareDocumentEnvironment{definition}{O{} O{}}{\slides@block{defcolor}{Definition}{#1}{#2}{6}}{\end{block}}
\DeclareDocumentEnvironment{example}{O{} O{}}{\slides@block{excolor}{Example}{#1}{#2}{6}}{\end{block}}
\DeclareDocumentEnvironment{remark}{O{} O{}}{\slides@block{gray!70!black}{Remark}{#1}{#2}{5}}{\end{block}}
\DeclareDocumentEnvironment{mlbox}{O{In machine learning}}{%
    \setbeamercolor{block title}{fg=white, bg=mlcolor}%
    \setbeamercolor{block body}{fg=black, bg=mlcolor!6}%
    \begin{block}{#1}}{\end{block}}
\makeatother

% Board pages.
%   \board{Title}{Task, shown in small print at the top}            one white page
%   \board[3]{Title}{Task}                                          three pages (the last two blank)
\newcounter{boardpage}
\newcommand{\boardhead}[1]{%
    \vskip0.3em
    \noindent{\color{boardink}\scriptsize
        \tikz[baseline=(b.base)]\node[draw=boardink, rounded corners=2pt, inner sep=1.5pt,
            font=\tiny\bfseries\sffamily] (b) {BOARD};\ \ #1}\par\vskip0.3em}
\newcommand{\board}[3][1]{%
    \begin{frame}[t]
        \boardhead{#2}
        \if\relax\detokenize{#3}\relax\else
            {\small #3\par}
        \fi
    \end{frame}%
    \setcounter{boardpage}{1}%
    \loop\ifnum\value{boardpage}<#1\relax
        \stepcounter{boardpage}%
        \begin{frame}[t]\boardhead{#2 (continued)}\end{frame}%
    \repeat}

% Practice pages: problems of the practice file, as white pages for solving them live.
% A deck loads the practice file's labels with \externaldocument[P1-]{../Lecture1/practice1} and
% \setpractice{1}. Then
%   \practice{pr:eig2}{\mandatory}{Problem text}    one page, headed "PRACTICE  Problem 10.1 (mandatory)"
%   \practice[2]{pr:eig-props}{}{Problem text}      two pages (the second blank)
%   \practice{pr:eig-gre}{$\star$}{Problem text}    the marks of the practice file
% The answers stay in LectureN/practiceN_answers.tex.
\newcommand{\mandatory}{\textcolor{red!70!black}{\textbf{mandatory}}}
\makeatletter
\newcommand{\setpractice}[1]{\def\slides@practice{#1}}
\newcommand{\practiceref}[1]{\ref*{P\slides@practice-#1}}
\newcommand{\practicehead}[2]{% label in the practice file, text after the problem number
    \vskip0.3em
    \noindent{\color{boardink}\scriptsize
        \tikz[baseline=(b.base)]\node[draw=defcolor, text=defcolor, rounded corners=2pt, inner sep=1.5pt,
            font=\tiny\bfseries\sffamily] (b) {PRACTICE};\ \ Practice~\slides@practice, Problem~\practiceref{#1}#2}%
    \par\vskip0.3em}
\newcommand{\practice}[4][1]{%
    \begin{frame}[t]
        \practicehead{#2}{\if\relax\detokenize{#3}\relax\else\ (#3)\fi}
        {\small #4\par}
    \end{frame}%
    \setcounter{boardpage}{1}%
    \loop\ifnum\value{boardpage}<#1\relax
        \stepcounter{boardpage}%
        \begin{frame}[t]\practicehead{#2}{\ (continued)}\end{frame}%
    \repeat}
\makeatother

% Divider page: \topic{\secref{sec:spectral}}{The spectral theorem} or \topic{\S2.1--2.2}{Norms}.
\newcommand{\topic}[2]{%
    \begin{frame}[c]
        \centering
        {\color{gray}\large #1}\par\vskip0.4em
        {\color{accent}\LARGE\bfseries #2}\par
    \end{frame}}

% Title page.
\setbeamercolor{title}{fg=accent}
\setbeamerfont{title}{series=\bfseries, size=\Large}
\setbeamercolor{subtitle}{fg=gray!60!black}
\setbeamerfont{subtitle}{size=\normalsize}
\setbeamertemplate{title page}{%
    \vfill
    {\usebeamercolor[fg]{title}\usebeamerfont{title}\inserttitle\par}
    \vskip0.6em
    {\usebeamercolor[fg]{subtitle}\usebeamerfont{subtitle}\insertsubtitle\par}
    \vskip1.2em
    {\small\color{gray}\insertauthor\par}
    \vfill}
\author{Mathematics for Machine Learning}
\date{}

% Two columns with a picture: \sidebyside{0.5}{left text}{right picture}.
\newcommand{\sidebyside}[3]{%
    \begin{columns}[T]
        \begin{column}{#1\linewidth}#2\end{column}
        \begin{column}{\dimexpr\linewidth-#1\linewidth-1em\relax}#3\end{column}
    \end{columns}}

% "Check" slide after a board computation.
\newcommand{\checkslide}[2]{%
    \begin{frame}{Check: #1}
        #2
    \end{frame}}

% Pictures.
\newcommand{\notesfig}[2][1]{\resizebox{#1\linewidth}{!}{\input{figures/notes/#2}}}
\newcommand{\fig}[2][1]{\includegraphics[width=#1\linewidth]{figures/#2}}

% aspectratio=1610 (16:10) matches the Samsung Galaxy Tab S9 FE+ (2560 x 1600) used for teaching.
%
% The decks have no overlays (\pause, \only): every slide is one PDF page, so what you write on a
% page while sharing the tablet screen stays on that page. Board pages (\board) are white pages with
% a small header, for calculations and proofs written live.

\usepackage[T1]{fontenc}
\usepackage{lmodern}
\usefonttheme[onlymath]{serif}
\usepackage{amsmath,amssymb}
\usepackage{booktabs}
\usepackage{tikz}
\usetikzlibrary{arrows.meta,calc,positioning,shapes.geometric,decorations.pathreplacing}
\usepackage{pgfplots}
\pgfplotsset{compat=1.18}
\usepgfplotslibrary{fillbetween}

% Numbers of the lecture notes. A deck says which notes it follows with \setnotes{2} and
% loads their labels with \externaldocument[L2-]{../Lecture2/lecture2} (the .aux of the notes
% must exist). Then \begin{theorem}[Spectral theorem][thm:spectral] prints "notes L2, 10.1" at
% the right of the block title, and \secref{sec:spectral} prints "§2.10".
\usepackage{xr-hyper}
\makeatletter
\newcommand{\setnotes}[1]{\def\slides@notes{#1}}
\setnotes{2}
\newcommand{\notesref}[1]{\ref*{L\slides@notes-#1}}
\newcommand{\secref}[1]{\S\slides@notes.\ref*{L\slides@notes-#1}}
\makeatother

% Macros of the lecture notes.
\newcommand{\norm}[1]{\lVert #1 \rVert}
\newcommand{\ip}[2]{\langle #1, #2 \rangle}
\newcommand{\py}[1]{\texttt{\detokenize{#1}}}
\newcommand{\E}{\mathbb{E}}
\newcommand{\indep}{\perp\!\!\!\perp}
\newcommand{\R}{\mathbb{R}}
\newcommand{\bx}{\mathbf{x}}
\newcommand{\by}{\mathbf{y}}
\newcommand{\bw}{\mathbf{w}}
\newcommand{\bv}{\mathbf{v}}
\newcommand{\bu}{\mathbf{u}}
\newcommand{\bq}{\mathbf{q}}
\newcommand{\bb}{\mathbf{b}}
\newcommand{\be}{\mathbf{e}}
\DeclareMathOperator{\Var}{Var}
\DeclareMathOperator{\Cov}{Cov}
\DeclareMathOperator{\Corr}{Corr}
\DeclareMathOperator{\KL}{KL}
\DeclareMathOperator{\im}{im}
\DeclareMathOperator{\rank}{rank}
\DeclareMathOperator{\Null}{null}
\DeclareMathOperator{\col}{col}
\DeclareMathOperator{\row}{row}
\DeclareMathOperator{\Span}{span}
\DeclareMathOperator{\tr}{tr}
\DeclareMathOperator{\diag}{diag}
\DeclareMathOperator{\Div}{div}
\DeclareMathOperator{\curl}{curl}
\DeclareMathOperator{\vol}{vol}
\DeclareMathOperator{\area}{area}
\DeclareMathOperator{\sign}{sign}
\DeclareMathOperator{\softmax}{softmax}
\DeclareMathOperator{\ReLU}{ReLU}
\DeclareMathOperator*{\argmin}{arg\,min}
\DeclareMathOperator*{\argmax}{arg\,max}

% Colors.
\definecolor{accent}{RGB}{31,78,140}      % titles, theorems
\definecolor{defcolor}{RGB}{24,110,92}    % definitions
\definecolor{excolor}{RGB}{150,82,20}     % examples
\definecolor{mlcolor}{RGB}{120,40,120}    % machine learning connections
\definecolor{boardink}{RGB}{140,140,140}  % header of the board pages

% Layout.
\setbeamertemplate{navigation symbols}{}
\setbeamersize{text margin left=0.9em, text margin right=0.9em}
\setbeamercolor{structure}{fg=accent}
\setbeamercolor{normal text}{fg=black, bg=white}
\setbeamercolor{frametitle}{fg=accent, bg=white}
\setbeamerfont{frametitle}{series=\bfseries, size=\large}
\setbeamertemplate{frametitle}{%
    \vskip0.45em\insertframetitle\par\vskip-0.25em
    {\color{accent!25}\rule{\linewidth}{0.5pt}}\vskip-0.2em}
\setbeamertemplate{itemize item}{\color{accent}\raisebox{0.15ex}{\scriptsize$\blacktriangleright$}}
\setbeamertemplate{itemize subitem}{\color{accent}\raisebox{0.2ex}{\tiny$\bullet$}}
\setbeamertemplate{enumerate item}{\color{accent}\insertenumlabel.}
\setbeamerfont{footline}{size=\tiny}
\setbeamertemplate{footline}{%
    \leavevmode\hbox to \paperwidth{\hskip0.9em{\color{gray}\insertshorttitle}\hfill
    {\color{gray}\insertframenumber}\hskip0.9em}\vskip0.45em}
\setbeamertemplate{blocks}[rounded][shadow=false]
\setbeamercolor{block title}{fg=white, bg=accent}
\setbeamercolor{block body}{fg=black, bg=accent!6}
\setbeamerfont{block title}{size=\normalsize, series=\bfseries}
\setlength{\abovedisplayskip}{4pt}
\setlength{\belowdisplayskip}{4pt}

% Blocks: \begin{theorem}[Spectral theorem][thm:spectral] ... \end{theorem}; both arguments are optional.
\makeatletter
\newcommand{\slides@block}[5]{% color, word, name, label in the notes, body tint
    \setbeamercolor{block title}{fg=white, bg=#1}%
    \setbeamercolor{block body}{fg=black, bg=#1!#5}%
    \begin{block}{#2\if\relax\detokenize{#3}\relax\else\ {\normalfont(#3)}\fi
        \if\relax\detokenize{#4}\relax\else\hfill{\normalfont\scriptsize notes~L\slides@notes, \notesref{#4}}\fi}}
\DeclareDocumentEnvironment{theorem}{O{} O{}}{\slides@block{accent}{Theorem}{#1}{#2}{6}}{\end{block}}
\DeclareDocumentEnvironment{proposition}{O{} O{}}{\slides@block{accent}{Proposition}{#1}{#2}{6}}{\end{block}}
\DeclareDocumentEnvironment{lemma}{O{} O{}}{\slides@block{accent}{Lemma}{#1}{#2}{6}}{\end{block}}
\DeclareDocumentEnvironment{corollary}{O{} O{}}{\slides@block{accent}{Corollary}{#1}{#2}{6}}{\end{block}}
\DeclareDocumentEnvironment{definition}{O{} O{}}{\slides@block{defcolor}{Definition}{#1}{#2}{6}}{\end{block}}
\DeclareDocumentEnvironment{example}{O{} O{}}{\slides@block{excolor}{Example}{#1}{#2}{6}}{\end{block}}
\DeclareDocumentEnvironment{remark}{O{} O{}}{\slides@block{gray!70!black}{Remark}{#1}{#2}{5}}{\end{block}}
\DeclareDocumentEnvironment{mlbox}{O{In machine learning}}{%
    \setbeamercolor{block title}{fg=white, bg=mlcolor}%
    \setbeamercolor{block body}{fg=black, bg=mlcolor!6}%
    \begin{block}{#1}}{\end{block}}
\makeatother

% Board pages.
%   \board{Title}{Task, shown in small print at the top}            one white page
%   \board[3]{Title}{Task}                                          three pages (the last two blank)
\newcounter{boardpage}
\newcommand{\boardhead}[1]{%
    \vskip0.3em
    \noindent{\color{boardink}\scriptsize
        \tikz[baseline=(b.base)]\node[draw=boardink, rounded corners=2pt, inner sep=1.5pt,
            font=\tiny\bfseries\sffamily] (b) {BOARD};\ \ #1}\par\vskip0.3em}
\newcommand{\board}[3][1]{%
    \begin{frame}[t]
        \boardhead{#2}
        \if\relax\detokenize{#3}\relax\else
            {\small #3\par}
        \fi
    \end{frame}%
    \setcounter{boardpage}{1}%
    \loop\ifnum\value{boardpage}<#1\relax
        \stepcounter{boardpage}%
        \begin{frame}[t]\boardhead{#2 (continued)}\end{frame}%
    \repeat}

% Practice pages: problems of the practice file, as white pages for solving them live.
% A deck loads the practice file's labels with \externaldocument[P1-]{../Lecture1/practice1} and
% \setpractice{1}. Then
%   \practice{pr:eig2}{\mandatory}{Problem text}    one page, headed "PRACTICE  Problem 10.1 (mandatory)"
%   \practice[2]{pr:eig-props}{}{Problem text}      two pages (the second blank)
%   \practice{pr:eig-gre}{$\star$}{Problem text}    the marks of the practice file
% The answers stay in LectureN/practiceN_answers.tex.
\newcommand{\mandatory}{\textcolor{red!70!black}{\textbf{mandatory}}}
\makeatletter
\newcommand{\setpractice}[1]{\def\slides@practice{#1}}
\newcommand{\practiceref}[1]{\ref*{P\slides@practice-#1}}
\newcommand{\practicehead}[2]{% label in the practice file, text after the problem number
    \vskip0.3em
    \noindent{\color{boardink}\scriptsize
        \tikz[baseline=(b.base)]\node[draw=defcolor, text=defcolor, rounded corners=2pt, inner sep=1.5pt,
            font=\tiny\bfseries\sffamily] (b) {PRACTICE};\ \ Practice~\slides@practice, Problem~\practiceref{#1}#2}%
    \par\vskip0.3em}
\newcommand{\practice}[4][1]{%
    \begin{frame}[t]
        \practicehead{#2}{\if\relax\detokenize{#3}\relax\else\ (#3)\fi}
        {\small #4\par}
    \end{frame}%
    \setcounter{boardpage}{1}%
    \loop\ifnum\value{boardpage}<#1\relax
        \stepcounter{boardpage}%
        \begin{frame}[t]\practicehead{#2}{\ (continued)}\end{frame}%
    \repeat}
\makeatother

% Divider page: \topic{\secref{sec:spectral}}{The spectral theorem} or \topic{\S2.1--2.2}{Norms}.
\newcommand{\topic}[2]{%
    \begin{frame}[c]
        \centering
        {\color{gray}\large #1}\par\vskip0.4em
        {\color{accent}\LARGE\bfseries #2}\par
    \end{frame}}

% Title page.
\setbeamercolor{title}{fg=accent}
\setbeamerfont{title}{series=\bfseries, size=\Large}
\setbeamercolor{subtitle}{fg=gray!60!black}
\setbeamerfont{subtitle}{size=\normalsize}
\setbeamertemplate{title page}{%
    \vfill
    {\usebeamercolor[fg]{title}\usebeamerfont{title}\inserttitle\par}
    \vskip0.6em
    {\usebeamercolor[fg]{subtitle}\usebeamerfont{subtitle}\insertsubtitle\par}
    \vskip1.2em
    {\small\color{gray}\insertauthor\par}
    \vfill}
\author{Mathematics for Machine Learning}
\date{}

% Two columns with a picture: \sidebyside{0.5}{left text}{right picture}.
\newcommand{\sidebyside}[3]{%
    \begin{columns}[T]
        \begin{column}{#1\linewidth}#2\end{column}
        \begin{column}{\dimexpr\linewidth-#1\linewidth-1em\relax}#3\end{column}
    \end{columns}}

% "Check" slide after a board computation.
\newcommand{\checkslide}[2]{%
    \begin{frame}{Check: #1}
        #2
    \end{frame}}

% Pictures.
\newcommand{\notesfig}[2][1]{\resizebox{#1\linewidth}{!}{\input{figures/notes/#2}}}
\newcommand{\fig}[2][1]{\includegraphics[width=#1\linewidth]{figures/#2}}

\externaldocument[L4-]{../Lecture4/lecture4}
\setnotes{4}

\title[Lecture 9 · Probability I--II]{Lecture 9. Probability I--II:\\ Random Variables, Joint Distributions and Gaussians}
\subtitle{Notes: Lecture 4, \S4.1--4.6}

\begin{document}

\begin{frame}[plain]
    \titlepage
\end{frame}

\begin{frame}{Today}
    \begin{columns}[T]
        \begin{column}{0.52\linewidth}
            \begin{enumerate}\small
                \item Probability spaces, Bayes' theorem, independence
                \item Random variables, distributions, sampling
                \item Expectation, variance, inequalities
                \item Joint distributions, covariance, initialization of networks
                \item Conditional expectation, graphical models, MGFs
                \item Random vectors and covariance matrices
                \item The multivariate Gaussian
            \end{enumerate}
        \end{column}
        \begin{column}{0.44\linewidth}
            \begin{mlbox}[Used in]
                \small
                \begin{itemize}
                    \item weight initialization and batch normalization (DL)
                    \item naive Bayes, probabilistic classifiers (ML)
                    \item variational autoencoders (DL)
                    \item Gaussian processes (Lecture~11)
                \end{itemize}
            \end{mlbox}
        \end{column}
    \end{columns}
\end{frame}

\begin{frame}{Probability is linear algebra and calculus, with weights}
    \small
    \begin{itemize}
        \item The expectation is a \textbf{linear functional}: $\E[aX + bY + c] = a\E[X] + b\E[Y] + c$.
        \item The variance is a \textbf{quadratic form}: $\Var(\bu^\top\mathbf{X}) = \bu^\top\Sigma\bu$.
        \item Covariance is an \textbf{inner product} on centered random variables; correlation is the cosine of an angle (Lecture~6).
        \item A covariance matrix is a \textbf{Gram matrix}: symmetric PSD, so the spectral theorem applies; its top eigenvector is the first principal component.
        \item Densities transform with \textbf{Jacobian determinants} (Lecture~7); Gaussian conditionals use the \textbf{Schur complement} (Lecture~1).
        \item $\E[Y \mid X]$ is an \textbf{orthogonal projection}: the best predictor of $Y$ from $X$.
    \end{itemize}
    Today we build the probability side; next time we let $N \to \infty$.
\end{frame}

% ------------------------------------------------------------------ Probability spaces
\topic{\secref{sec:probability}}{Probability spaces}

\begin{frame}{Probability spaces}
    \begin{definition}[Probability space][def:probability]
        \small A sample space $\Omega$ and a probability $P$ on events $A \subseteq \Omega$: \ $P(A) \geq 0$, \ $P(\Omega) = 1$, \ $P(A_1 \cup A_2 \cup \cdots) = \sum_iP(A_i)$ for disjoint $A_i$.
    \end{definition}
    \begin{proposition}[][prop:probability-rules]
        \small $P(A^c) = 1 - P(A)$; \ $A \subseteq B \Rightarrow P(A) \leq P(B)$; \ $P(A \cup B) = P(A) + P(B) - P(A \cap B)$; \ union bound $P(\bigcup_iA_i) \leq \sum_iP(A_i)$.
    \end{proposition}
    \begin{example}[][ex:dice]
        \small Two fair dice: $P(\text{sum} = 7) = \frac{6}{36}$; \ $P(\text{at least one six}) = 1 - \frac{25}{36} = \frac{11}{36}$.
    \end{example}
    \small For uncountable $\Omega$ not every subset can be an event: $\sigma$-algebras and measure theory, in Appendix~B. We ignore this: intervals, balls and their unions are always events.
\end{frame}

\begin{frame}{Conditional probability and Bayes' theorem}
    \sidebyside{0.56}{%
        \small
        $P(A \mid B) = \frac{P(A \cap B)}{P(B)}$ \ for $P(B) > 0$.
        \begin{theorem}[Total probability, Bayes][thm:bayes]
            \small For a partition $B_1, \ldots, B_n$:
            \[P(A) = \sum_iP(A \mid B_i)P(B_i),\]
            \[P(B_j \mid A) = \frac{P(A \mid B_j)P(B_j)}{\sum_iP(A \mid B_i)P(B_i)}.\]
        \end{theorem}
        Independence: $P(A \cap B) = P(A)P(B)$; for several events, for every sub-collection.
    }{%
        \centering\fig[1]{s09_bayes}
    }
\end{frame}

\board{Example: a medical test}{A disease affects $1\%$ of the population. The test is positive for $99\%$ of the sick, and for $5\%$ of the healthy. What is $P(\text{sick} \mid \text{positive})$? Count among $10\,000$ people.}

\checkslide{the base rate fallacy}{
    \small
    \textbf{Bayes:} $P(D) = 0.01$, $P(+ \mid D) = 0.99$, $P(+ \mid D^c) = 0.05$. Total probability:
    \[P(+) = 0.99 \cdot 0.01 + 0.05 \cdot 0.99 = 0.0099 + 0.0495 = 0.0594, \qquad P(D \mid +) = \frac{0.0099}{0.0594} = \frac{1}{6} \approx 0.167.\]
    \textbf{Counting} $10\,000$ people:
    \begin{center}
        \renewcommand{\arraystretch}{1.2}
        \begin{tabular}{@{}lccc@{}}
            \toprule
            & positive & negative & total \\
            \midrule
            sick & $0.99 \cdot 100 = 99$ & $1$ & $100$ \\
            healthy & $0.05 \cdot 9900 = 495$ & $9405$ & $9900$ \\
            \midrule
            total & $594$ & $9406$ & $10\,000$ \\
            \bottomrule
        \end{tabular}
    \end{center}
    Among the $594$ positives, $99$ are sick: $\frac{99}{594} = \frac{1}{6}$. A second independent positive test: prior odds $\frac{1}{5}$, times the likelihood ratio $\frac{0.99}{0.05} = 19.8$, gives odds $3.96$, i.e., $P(D) \approx 0.80$.
    \\[0.3em]
    The test is accurate, but the disease is rare: ignoring the prior $P(D)$ is the \textbf{base rate fallacy}.
    \begin{mlbox}
        \small Classifiers for rare classes (fraud, rare diseases) produce many false alarms even with a small false positive rate: look at precision, not only accuracy (self-study \S4.11).
    \end{mlbox}
}

% ------------------------------------------------------------------ Random variables
\topic{\secref{sec:random-variables}}{Random variables and their distributions}

\begin{frame}{Random variables}
    \small
    A random variable is a function $X: \Omega \to \R$; its CDF is $F(x) = P(X \leq x)$.
    \begin{definition}[][def:pmf-pdf]
        \small \textbf{Discrete}: pmf $p(x) = P(X = x)$, $\sum_ip(x_i) = 1$. \ \textbf{Continuous}: density $p \geq 0$ with $P(a \leq X \leq b) = \int_a^bp(x)\,dx$.
    \end{definition}
    \sidebyside{0.48}{%
        \small
        \begin{itemize}
            \item Bernoulli$(p)$, Binomial$(n, p)$: $\binom{n}{k}p^k(1 - p)^{n-k}$
            \item Geometric$(p)$: $(1 - p)^{k-1}p$
            \item Poisson$(\lambda)$: $e^{-\lambda}\frac{\lambda^k}{k!}$
            \item Categorical$(\mathbf{p})$: softmax outputs
            \item Uniform$(a, b)$; Exponential$(\lambda)$: $\lambda e^{-\lambda x}$, memoryless
            \item $\mathcal{N}(\mu, \sigma^2)$: $\frac{1}{\sqrt{2\pi\sigma^2}}e^{-(x - \mu)^2/2\sigma^2}$
        \end{itemize}
    }{%
        \centering\notesfig[1]{L4-distributions}
    }
\end{frame}

\begin{frame}{Transformations and sampling}
    \begin{proposition}[Change of variables][prop:change-of-variables-1d]
        \small $g$ differentiable, strictly monotone: \ $Y = g(X)$ has the density $p_Y(y) = p_X(x)/|g'(x)|$, $x = g^{-1}(y)$.
    \end{proposition}
    \small $Z = \frac{X - \mu}{\sigma} \sim \mathcal{N}(0, 1)$; \ $P(|X - \mu| \leq k\sigma) = 0.683, 0.954, 0.997$ ($k = 1, 2, 3$).
    \begin{theorem}[Inverse transform sampling][thm:inverse-transform]
        \small With $F^{-1}(u) = \min\{x : F(x) \geq u\}$ and $U \sim$ Uniform$(0, 1)$, $X = F^{-1}(U)$ has the CDF $F$.
    \end{theorem}
    \centering\fig[0.56]{s09_inverse_transform}
\end{frame}

\board{Proof of inverse transform sampling}{Show $F^{-1}(u) \leq x \iff u \leq F(x)$; then $P(F^{-1}(U) \leq x) = P(U \leq F(x)) = F(x)$. Example: sample Exponential$(\lambda)$, and a discrete distribution with values $1, 2, 3$ and probabilities $0.2, 0.5, 0.3$.}

\checkslide{distributions and sampling with numbers}{
    \small
    \textbf{Binomial$(10, 0.3)$:} $P(X = 3) = \binom{10}{3}0.3^3\,0.7^7 = 120 \cdot 0.027 \cdot 0.0824 \approx 0.267$; \ $\E X = 3$, $\Var X = 2.1$.
    \\[0.2em]
    \textbf{Exponential, memoryless:} $P(X > s + t \mid X > s) = \frac{e^{-\lambda(s + t)}}{e^{-\lambda s}} = e^{-\lambda t} = P(X > t)$.
    \\[0.2em]
    \textbf{Standardization:} $X \sim \mathcal{N}(100, 15^2)$: $P(X > 130) = P(Z > \frac{130 - 100}{15}) = P(Z > 2) \approx 0.023$.
    \\[0.4em]
    \textbf{Inverse transform, exponential:} $F(x) = 1 - e^{-\lambda x}$; solve $u = 1 - e^{-\lambda x}$: \ $F^{-1}(u) = -\frac{\log(1 - u)}{\lambda}$. \ $u = 0.5$ gives the median $\frac{\log 2}{\lambda} \approx \frac{0.693}{\lambda}$; \ $u = 0.9$ gives $\frac{2.303}{\lambda}$.
    \\[0.4em]
    \textbf{Inverse transform, discrete:} values $1, 2, 3$ with probabilities $0.2, 0.5, 0.3$; the CDF jumps to $0.2, 0.7, 1$:
    \[F^{-1}(u) = \min\{x : F(x) \geq u\} = \begin{cases} 1, & 0 < u \leq 0.2, \\ 2, & 0.2 < u \leq 0.7, \\ 3, & 0.7 < u < 1. \end{cases} \qquad u = 0.65 \mapsto 2, \quad u = 0.93 \mapsto 3.\]
    The intervals have lengths $0.2, 0.5, 0.3$: exactly the probabilities. (This is \py{np.searchsorted(cdf, u)}.)
}

% ------------------------------------------------------------------ Expectation
\topic{\secref{sec:expectation}}{Expectation and variance}

\begin{frame}{Expectation and variance}
    \small
    \begin{definition}[][def:expectation]
        \small $\E[X] = \sum_ix_ip(x_i)$ or $\int xp(x)\,dx$ (if absolutely convergent); \ $\E[g(X)] = \int g(x)p(x)\,dx$.
    \end{definition}
    Linear: $\E[aX + bY + c] = a\E[X] + b\E[Y] + c$. \ $\Var(X) = \E[(X - \mu)^2] = \E[X^2] - \E[X]^2$, \ $\Var(aX + b) = a^2\Var(X)$.
    \begin{columns}[T]
        \begin{column}{0.5\linewidth}
            \centering\small
            \renewcommand{\arraystretch}{1.15}
            \begin{tabular}{@{}lll@{}}
                \toprule
                & $\E[X]$ & $\Var(X)$ \\
                \midrule
                Bernoulli$(p)$ & $p$ & $p(1 - p)$ \\
                Binomial$(n, p)$ & $np$ & $np(1 - p)$ \\
                Geometric$(p)$ & $\frac{1}{p}$ & $\frac{1 - p}{p^2}$ \\
                Poisson$(\lambda)$ & $\lambda$ & $\lambda$ \\
                Uniform$(a, b)$ & $\frac{a + b}{2}$ & $\frac{(b - a)^2}{12}$ \\
                Exponential$(\lambda)$ & $\frac{1}{\lambda}$ & $\frac{1}{\lambda^2}$ \\
                $\mathcal{N}(\mu, \sigma^2)$ & $\mu$ & $\sigma^2$ \\
                \bottomrule
            \end{tabular}
        \end{column}
        \begin{column}{0.46\linewidth}
            \begin{example}[Tail sums][ex:tail-sum]
                \small For $X \in \{0, 1, 2, \ldots\}$: \ $\E[X] = \sum_{k \geq 0}P(X > k)$, since $X = \sum_{k \geq 0}\mathbf{1}_{\{X > k\}}$. \\
                Geometric: $\sum_k(1 - p)^k = \frac{1}{p}$. \\
                For $X \geq 0$ with a density: $\E[X] = \int_0^\infty P(X > t)\,dt$.
            \end{example}
        \end{column}
    \end{columns}
\end{frame}

\board{Compute means and variances}{Fair die: $\E[X]$, $\E[X^2]$, $\Var(X)$. \quad Exponential$(\lambda)$: $\E[X]$ with the tail integral. \quad Binomial$(n, p)$: as a sum of $n$ Bernoulli variables.}

\checkslide{means and variances}{
    \small
    \textbf{Fair die:} $\E X = \frac{1 + 2 + \cdots + 6}{6} = \frac{21}{6} = 3.5$, \ $\E X^2 = \frac{1 + 4 + 9 + 16 + 25 + 36}{6} = \frac{91}{6}$,
    \[\Var X = \tfrac{91}{6} - 3.5^2 = \tfrac{91}{6} - \tfrac{49}{4} = \tfrac{182 - 147}{12} = \tfrac{35}{12} \approx 2.92, \qquad \sigma \approx 1.71.\]
    \textbf{Exponential$(\lambda)$:} $P(X > t) = e^{-\lambda t}$, so $\E X = \int_0^\infty e^{-\lambda t}\,dt = \frac{1}{\lambda}$ (no integration by parts).
    \\[0.3em]
    \textbf{Geometric$(p)$:} $P(X > k) = (1 - p)^k$ (the first $k$ trials fail), so $\E X = \sum_{k \geq 0}(1 - p)^k = \frac{1}{1 - (1 - p)} = \frac{1}{p}$. A fair die needs $6$ throws on average for a six.
    \\[0.3em]
    \textbf{Binomial$(n, p)$} $= X_1 + \cdots + X_n$ with independent Bernoulli$(p)$: $\E X_i = p$, $\E X_i^2 = p$, $\Var X_i = p - p^2$. Linearity, and variances of independent variables add:
    \[\E X = np, \qquad \Var X = np(1 - p). \qquad n = 100, p = \tfrac{1}{2}: \ \E X = 50, \ \sigma = 5.\]
}

\begin{frame}{Inequalities}
    \begin{proposition}[Markov and Chebyshev][prop:markov]
        \small $X \geq 0$: \ $P(X \geq a) \leq \frac{\E[X]}{a}$. \qquad Mean $\mu$, variance $\sigma^2$: \ $P(|X - \mu| \geq k) \leq \frac{\sigma^2}{k^2}$.
    \end{proposition}
    \begin{theorem}[Jensen][thm:jensen]
        \small $f$ convex: $f(\E[X]) \leq \E[f(X)]$; \ $f$ concave: $\geq$. \quad E.g.\ $\E[X]^2 \leq \E[X^2]$, \ $\E[\log X] \leq \log\E[X]$.
    \end{theorem}
    \begin{example}[The Chernoff bound][ex:chernoff]
        \small $P(X \geq a) = P(e^{tX} \geq e^{ta}) \leq e^{-ta}\E[e^{tX}]$, then choose the best $t > 0$. \\
        $S \sim$ Binomial$(100, \frac{1}{2})$: \ Chernoff $P(S \geq 60) \leq 0.134$, \ Chebyshev $\leq 0.25$, \ exact $0.028$. Chernoff bounds decay exponentially in $n$, Chebyshev's like $\frac{1}{n}$.
    \end{example}
\end{frame}

\board{Proofs: Markov, Chebyshev, Jensen}{Markov: $a\mathbf{1}_{\{X \geq a\}} \leq X$. \quad Chebyshev: Markov for $(X - \mu)^2$. \quad Jensen: a convex $f$ lies above its tangent line at $\E[X]$ (Lecture~8).}

\checkslide{three bounds for $P(S \geq 60)$, $S \sim$ Binomial$(100, \frac{1}{2})$}{
    \small
    $\E S = 50$, $\Var S = 25$.
    \\[0.3em]
    \textbf{Markov} ($S \geq 0$): $P(S \geq 60) \leq \frac{50}{60} \approx 0.83$. \hfill (uses only the mean)
    \\[0.2em]
    \textbf{Chebyshev:} $P(S \geq 60) \leq P(|S - 50| \geq 10) \leq \frac{25}{10^2} = 0.25$. \hfill (uses the variance)
    \\[0.2em]
    \textbf{Chernoff:} $\E[e^{tS}] = \prod_i\E[e^{tX_i}] = \big(\frac{1 + e^t}{2}\big)^{100}$, so
    \[P(S \geq 60) \leq \min_{t > 0}e^{-60t}\Big(\frac{1 + e^t}{2}\Big)^{100}.\]
    Minimize the log: $-60 + 100\frac{e^t}{1 + e^t} = 0 \iff e^t = \frac{60}{40} = 1.5$. Then
    \[P(S \geq 60) \leq \Big(\frac{2}{3}\Big)^{60}\Big(\frac{2.5}{2}\Big)^{100} = \exp\big(60\log\tfrac{2}{3} + 100\log 1.25\big) = e^{-24.33 + 22.31} \approx e^{-2.01} \approx 0.134.\]
    \textbf{Exact:} $\sum_{k \geq 60}\binom{100}{k}2^{-100} \approx 0.028$. \ With $n = 1000$ and $S \geq 600$, Chebyshev gives $0.025$ but Chernoff gives $\approx 2 \cdot 10^{-9}$.
}

% ------------------------------------------------------------------ Joint
\topic{\secref{sec:joint}}{Joint distributions and independence}

\begin{frame}{Joint, marginal and conditional distributions}
    \small
    \begin{definition}[][def:joint]
        \small Joint pmf $p(x, y)$ or density; \ marginal $p_X(x) = \sum_yp(x, y)$ (sum rule); \ conditional $p(y \mid x) = \frac{p(x, y)}{p_X(x)}$, i.e., $p(x, y) = p(y \mid x)p_X(x)$ (product rule).
    \end{definition}
    Independent: $p(x, y) = p_X(x)p_Y(y)$; then $\E[g(X)h(Y)] = \E[g(X)]\E[h(Y)]$.
    \begin{columns}[T]
        \begin{column}{0.42\linewidth}
            \centering
            \renewcommand{\arraystretch}{1.15}
            \begin{tabular}{c|cc|c}
                & $Y = 0$ & $Y = 1$ & $p_X$ \\
                \hline
                $X = 0$ & $0.3$ & $0.2$ & $0.5$ \\
                $X = 1$ & $0.1$ & $0.4$ & $0.5$ \\
                \hline
                $p_Y$ & $0.4$ & $0.6$ &
            \end{tabular}
        \end{column}
        \begin{column}{0.54\linewidth}
            \small
            $P(Y = 1 \mid X = 1) = 0.8$, \ $P(Y = 1 \mid X = 0) = 0.4$. \\
            $p(1, 1) = 0.4 \neq 0.5 \cdot 0.6$: not independent. \\
            $\Cov(X, Y) = 0.4 - 0.3 = 0.1$, \ $\Corr \approx 0.41$.
        \end{column}
    \end{columns}
    \vskip0.5em
    \textbf{Sums of independent variables} have the convolution density $p_{X+Y}(z) = \int p_X(x)p_Y(z - x)\,dx$: two Uniform$(0, 1)$ give a triangle; Poisson$(\lambda)$ $+$ Poisson$(\mu)$ $=$ Poisson$(\lambda + \mu)$.
\end{frame}

\begin{frame}{Covariance and correlation}
    \small
    \begin{definition}[][def:covariance]
        \small $\Cov(X, Y) = \E[(X - \mu_X)(Y - \mu_Y)] = \E[XY] - \E[X]\E[Y]$, \ $\Corr(X, Y) = \frac{\Cov(X, Y)}{\sigma_X\sigma_Y}$.
    \end{definition}
    \begin{proposition}[][prop:variance-sum]
        \small $\Var(X + Y) = \Var X + \Var Y + 2\Cov(X, Y)$; for pairwise uncorrelated variables the variances add.
    \end{proposition}
    \begin{theorem}[][thm:correlation]
        \small $|\Corr(X, Y)| \leq 1$, with equality iff $Y = aX + b$ ($a \neq 0$) with probability $1$.
    \end{theorem}
    \centering\fig[0.74]{s09_correlation}
    \\[0.1em]
    \raggedright Uncorrelated $\not\Rightarrow$ independent: correlation only sees \emph{linear} dependence.
\end{frame}

\board{Proof: $|\Corr(X, Y)| \leq 1$}{$\ip{X}{Y} = \E[XY]$ is an inner product on centered random variables (up to $X = 0$ with probability $1$). Apply Cauchy--Schwarz (Lecture~6). Equality case?}

\begin{frame}{Worked examples: joint distributions and sums}
    \small
    \textbf{The $2 \times 2$ table} ($p(0, 0) = 0.3$, $p(0, 1) = 0.2$, $p(1, 0) = 0.1$, $p(1, 1) = 0.4$):
    \begin{itemize}
        \item marginals: $P(X = 1) = 0.1 + 0.4 = 0.5$, \ $P(Y = 1) = 0.2 + 0.4 = 0.6$;
        \item $\E[XY] = 1 \cdot 1 \cdot 0.4 = 0.4$, \ $\Cov = 0.4 - 0.5 \cdot 0.6 = 0.1$;
        \item $\Var X = 0.5 \cdot 0.5 = 0.25$, $\Var Y = 0.6 \cdot 0.4 = 0.24$, \ $\Corr = \frac{0.1}{\sqrt{0.25 \cdot 0.24}} = \frac{0.1}{0.245} \approx 0.41$.
    \end{itemize}
    \textbf{Two uniforms.} $X, Y \sim$ Uniform$(0, 1)$ independent; the convolution is $p_{X+Y}(z) = \int_0^1\mathbf{1}_{\{0 \leq z - x \leq 1\}}\,dx$, the length of $[\max(0, z - 1), \min(1, z)]$:
    \[p_{X+Y}(z) = z \ \ (0 \leq z \leq 1), \qquad 2 - z \ \ (1 \leq z \leq 2), \qquad P(X + Y \leq \tfrac{1}{2}) = \int_0^{1/2}z\,dz = \tfrac{1}{8}.\]
    \textbf{Two Poissons.} Clicks from search $\sim$ Poisson$(3)$, from ads $\sim$ Poisson$(2)$, independent:
    \[P(X + Y = n) = \sum_{k=0}^ne^{-3}\tfrac{3^k}{k!}e^{-2}\tfrac{2^{n-k}}{(n - k)!} = \tfrac{e^{-5}}{n!}\sum_k\tbinom{n}{k}3^k2^{n-k} = e^{-5}\tfrac{5^n}{n!}:\]
    Poisson$(5)$ (binomial theorem). E.g.\ $P(\text{no clicks in a minute}) = e^{-5} \approx 0.0067$.
\end{frame}

\begin{frame}{The variance of a layer, and initialization}
    \begin{proposition}[Variance of a layer][prop:layer-variance]
        \small $z = W_1x_1 + \cdots + W_nx_n$, weights independent of the $x_j$, mean $0$, variance $v$, $\E[x_j^2] = q$: \ $\E[z] = 0$, \ $\Var(z) = nvq$. If the $W_j$ are symmetric, $\E[\ReLU(z)^2] = \frac{1}{2}\Var(z)$.
    \end{proposition}
    \sidebyside{0.5}{%
        \small
        Deep ReLU network: $q_L = q_0\prod_l\frac{n_{l-1}v_l}{2}$.
        \begin{itemize}
            \item $v = \frac{2}{n_{\text{in}}}$ (\textbf{He}): signals keep their size;
            \item $v = \frac{1}{n}$: halved per layer, $2^{-50} \approx 10^{-15}$;
            \item $v = 1$, $n = 512$: $\times 256$ per layer.
        \end{itemize}
        $\tanh$: forward $nv = 1$, backward $n_{\text{out}}v = 1$: \textbf{Xavier} $v = \frac{2}{n_{\text{in}} + n_{\text{out}}}$.
    }{%
        \centering\fig[1]{s09_initialization}
    }
\end{frame}

\board{Proof of $\Var(z) = nvq$}{$\E[W_jx_j] = \E[W_j]\E[x_j] = 0$; \ $\E[W_jx_jW_kx_k] = \E[W_jW_k]\E[x_jx_k]$, which is $0$ for $j \neq k$ and $vq$ for $j = k$. \quad ReLU: $z$ and $-z$ have the same distribution, and $\ReLU(z)^2 + \ReLU(-z)^2 = z^2$.}

\checkslide{signal sizes through a ReLU network}{
    \small
    \textbf{One layer.} $\Var(z) = \E[z^2] = \sum_{j,k}\E[W_jW_k]\E[x_jx_k] = \sum_j v\,q = nvq$; \ then $\E[\ReLU(z)^2] = \frac{1}{2}\E[z^2]$, since $\E[\ReLU(z)^2] = \E[\ReLU(-z)^2]$ and they add up to $\E[z^2]$.
    \\[0.4em]
    \textbf{Per layer} the second moment is multiplied by $\frac{1}{2}nv$. With $n = 512$ units per layer:
    \begin{center}
        \renewcommand{\arraystretch}{1.25}
        \begin{tabular}{@{}lccl@{}}
            \toprule
            weight variance $v$ & factor $\frac{1}{2}nv$ & after $50$ layers & \\
            \midrule
            $1$ (standard normal) & $256$ & $256^{50} \approx 10^{120}$ & explodes \\
            $\frac{1}{512}$ & $\frac{1}{2}$ & $2^{-50} \approx 9 \cdot 10^{-16}$ & vanishes \\
            $\frac{2}{512}$ (He) & $1$ & $1$ & stable \\
            \bottomrule
        \end{tabular}
    \end{center}
    \textbf{Uniform weights} on $[-b, b]$ have variance $\frac{b^2}{3}$; He: $\frac{b^2}{3} = \frac{2}{n}$, $b = \sqrt{6/n} \approx 0.108$ for $n = 512$.
    \\[0.3em]
    \textbf{Xavier} for $\tanh$, $n_{\text{in}} = 512$, $n_{\text{out}} = 256$: $v = \frac{2}{768} \approx 0.0026$, a compromise between $\frac{1}{512}$ and $\frac{1}{256}$.
}

\begin{frame}{Conditional expectation}
    \begin{definition}[][def:conditional-expectation]
        \small $m(x) = \E[Y \mid X = x] = \sum_yy\,p(y \mid x)$; \ the random variable $\E[Y \mid X] = m(X)$.
    \end{definition}
    \begin{theorem}[][thm:conditional-expectation]
        \small
        \begin{enumerate}
            \item Tower property: $\E\big[\E[Y \mid X]\big] = \E[Y]$.
            \item Best prediction: $\E[(Y - g(X))^2] = \E[(Y - m(X))^2] + \E[(m(X) - g(X))^2]$.
        \end{enumerate}
    \end{theorem}
    \small
    $m(X)$ is the orthogonal projection of $Y$ onto the functions of $X$ (Pythagoras!). The regression function that every regression model tries to learn is $\E[Y \mid X = x]$.
    \\[0.3em]
    \textbf{Law of total variance:} \ $\Var(Y) = \E[\Var(Y \mid X)] + \Var(\E[Y \mid X])$: within-group variance plus variance of the group means.
\end{frame}

\board[2]{Examples of conditional expectation}{(1) $X_1, \ldots, X_n$ i.i.d., $S = \sum_iX_i$: show $\E[X_1 \mid S] = S/n$. \\ (2) Customers $N \sim$ Poisson$(\lambda)$, each spends $Y_i$ (i.i.d., mean $\mu$, variance $\sigma^2$, independent of $N$). Find the mean and variance of the revenue $T = Y_1 + \cdots + Y_N$.}

\checkslide{conditional expectations}{
    \small
    \textbf{(1) Symmetry.} The pairs $(X_i, S)$ all have the same distribution, so $\E[X_i \mid S] = g(S)$ for the same $g$; summing, $n\,g(S) = \E[S \mid S] = S$. \ Two dice with sum $8$: the outcomes $(2, 6), (3, 5), (4, 4), (5, 3), (6, 2)$ are equally likely, and the first die has mean $\frac{2 + 3 + 4 + 5 + 6}{5} = 4 = \frac{8}{2}$ \checkmark.
    \\[0.4em]
    \textbf{(2) A random sum.} Given $N = n$: $T$ is a sum of $n$ independent terms, so $\E[T \mid N] = N\mu$, $\Var(T \mid N) = N\sigma^2$. Tower property and total variance ($\E N = \Var N = \lambda$):
    \[\E[T] = \E[N\mu] = \lambda\mu, \qquad \Var(T) = \E[N\sigma^2] + \Var(N\mu) = \lambda\sigma^2 + \mu^2\lambda.\]
    \textbf{Numbers:} $\lambda = 20$ customers per day, each spends $\mu = 30$ with $\sigma = 10$:
    \[\E T = 600, \qquad \Var T = 20 \cdot 100 + 900 \cdot 20 = 20\,000, \qquad \mathrm{sd}(T) \approx 141.\]
    With a fixed $N = 20$, the sd would be $\sqrt{20} \cdot 10 \approx 45$: most of the uncertainty comes from the number of customers.
}

\begin{frame}{Conditional independence and graphical models}
    \small
    $X \indep Y \mid Z$ if $p(x, y \mid z) = p(x \mid z)p(y \mid z)$.
    \begin{proposition}[][prop:three-graphs]
        \small \textbf{Chain} $X \to Z \to Y$ and \textbf{fork} $X \leftarrow Z \to Y$: $X \indep Y \mid Z$. \ \textbf{Collider} $X \to Z \leftarrow Y$: $X \indep Y$, but in general not given $Z$ (explaining away).
    \end{proposition}
    \begin{example}[Naive Bayes][ex:naive-bayes]
        \small $p(\bx, y) = p(y)\prod_jp(x_j \mid y)$, so $p(y \mid \bx) \propto p(y)\prod_jp(x_j \mid y)$. Spam: $P(\text{spam}) = 0.4$; ``free'' in $30\%$ of spam, $2\%$ of other; ``meeting'' in $1\%$ and $10\%$. ``free'' without ``meeting'':
        \[\frac{P(\text{spam} \mid \bx)}{P(\text{other} \mid \bx)} = \frac{0.4}{0.6}\cdot\frac{0.3}{0.02}\cdot\frac{0.99}{0.9} = 11, \qquad P(\text{spam} \mid \bx) = \tfrac{11}{12}.\]
    \end{example}
    The log-odds of naive Bayes are affine in $\bx$, as in logistic regression: a generative and a discriminative model of the same form.
\end{frame}

\begin{frame}{Moment generating functions}
    \begin{definition}[][def:mgf]
        \small $M_X(t) = \E[e^{tX}]$; \ $K_X(t) = \log M_X(t)$ (cumulant generating function).
    \end{definition}
    \begin{proposition}[][prop:mgf]
        \small $M_{aX + b}(t) = e^{bt}M_X(at)$; \ independent $X$, $Y$: $M_{X+Y} = M_XM_Y$, $K_{X+Y} = K_X + K_Y$; \ $\E[X^k] = M_X^{(k)}(0)$, $K_X'(0) = \E[X]$, $K_X''(0) = \Var(X)$.
    \end{proposition}
    \small
    \textbf{Poisson$(\lambda)$:} $M(t) = \sum_ke^{tk}e^{-\lambda}\frac{\lambda^k}{k!} = e^{-\lambda}\sum_k\frac{(\lambda e^t)^k}{k!} = e^{\lambda(e^t - 1)}$, \ $K(t) = \lambda(e^t - 1)$: \ $K'(0) = \lambda = \E X$, $K''(0) = \lambda = \Var X$. \ Sums: $K_{X+Y} = (\lambda + \mu)(e^t - 1)$, so Poisson $+$ Poisson $=$ Poisson.
    \\[0.3em]
    \textbf{Bernoulli$(p)$:} $M(t) = 1 - p + pe^t$; \ Binomial$(n, p)$: $(1 - p + pe^t)^n$.
    \\[0.3em]
    \textbf{$\mathcal{N}(\mu, \sigma^2)$:} $M(t) = e^{\mu t + \sigma^2t^2/2}$, $K(t) = \mu t + \frac{\sigma^2t^2}{2}$: all cumulants beyond the second vanish.
    \\[0.3em]
    Products of MGFs make sums of independent variables easy: the Chernoff bound $P(S \geq a) \leq \min_te^{-ta}\prod_iM_{X_i}(t)$, and the central limit theorem (next lecture).
\end{frame}

% ------------------------------------------------------------------ Random vectors
\topic{\secref{sec:random-vectors}}{Random vectors and covariance matrices}

\begin{frame}{Covariance matrices}
    \begin{definition}[][def:covariance-matrix]
        \small $\boldsymbol{\mu} = \E[\mathbf{X}]$, \ $\Sigma = \E[(\mathbf{X} - \boldsymbol{\mu})(\mathbf{X} - \boldsymbol{\mu})^\top]$, \ $\Sigma_{ij} = \Cov(X_i, X_j)$.
    \end{definition}
    \begin{proposition}[][prop:covariance-matrix]
        \small $\Sigma = \E[\mathbf{X}\mathbf{X}^\top] - \boldsymbol{\mu}\boldsymbol{\mu}^\top$; \ $\Cov(A\mathbf{X} + \mathbf{b}) = A\Sigma A^\top$; \ $\Var(\bu^\top\mathbf{X}) = \bu^\top\Sigma\bu$, so $\Sigma$ is symmetric PSD.
    \end{proposition}
    \begin{example}[][ex:covariance-matrix]
        \small $\Sigma = \begin{pmatrix} 4 & 2 \\ 2 & 3 \end{pmatrix}$: $\Var(X_1 + X_2) = 11$, $\Var(X_1 - X_2) = 3$, $\Corr \approx 0.58$; eigenvalues $\frac{7 \pm \sqrt{17}}{2} \approx 5.56, 1.44$: the first principal component carries $5.56$ of the total variance $\tr\Sigma = 7$.
    \end{example}
    \begin{definition}[Whitening, Mahalanobis][def:whitening]
        \small $\mathbf{Z} = \Sigma^{-1/2}(\mathbf{X} - \boldsymbol{\mu})$ has $\Cov(\mathbf{Z}) = I$; \ $d_\Sigma(\bx, \boldsymbol{\mu}) = \norm{\Sigma^{-1/2}(\bx - \boldsymbol{\mu})}$.
    \end{definition}
\end{frame}

\board{Prove $\Cov(A\mathbf{X} + \mathbf{b}) = A\Sigma A^\top$}{Then: $\Var(\bu^\top\mathbf{X}) = \bu^\top\Sigma\bu \geq 0$; the direction of largest variance is the top eigenvector (Rayleigh, Lecture~6); whitening gives $\Sigma^{-1/2}\Sigma\Sigma^{-1/2} = I$. In tensor language: $\Sigma = \E[(\mathbf{X} - \boldsymbol{\mu}) \otimes (\mathbf{X} - \boldsymbol{\mu})]$, $\Cov(A\mathbf{X})_{ab} = A_{ai}A_{bj}\Sigma_{ij}$.}

\checkslide{covariance matrices with numbers}{
    \small
    \textbf{Proof.} $A\mathbf{X} + \mathbf{b} - (A\boldsymbol{\mu} + \mathbf{b}) = A(\mathbf{X} - \boldsymbol{\mu})$, so $\Cov = \E[A(\mathbf{X} - \boldsymbol{\mu})(\mathbf{X} - \boldsymbol{\mu})^\top A^\top] = A\Sigma A^\top$ (linearity).
    \\[0.4em]
    \textbf{Example} $\Sigma = \begin{pmatrix} 4 & 2 \\ 2 & 3 \end{pmatrix}$:
    \[\Var(X_1 + X_2) = \begin{pmatrix} 1 & 1 \end{pmatrix}\Sigma\begin{pmatrix} 1 \\ 1 \end{pmatrix} = 4 + 2 + 2 + 3 = 11, \qquad \Var(X_1 - X_2) = 4 - 2 - 2 + 3 = 3,\]
    \[\Corr(X_1, X_2) = \frac{2}{\sqrt{4 \cdot 3}} \approx 0.58; \qquad \tr\Sigma = 7, \ \det\Sigma = 8: \ \lambda = \frac{7 \pm \sqrt{49 - 32}}{2} \approx 5.56, \ 1.44.\]
    The first principal direction carries $\frac{5.56}{7} \approx 79\%$ of the total variance.
    \\[0.4em]
    \textbf{Joint covariance} of $(X_1 + X_2, X_1 - X_2)$: $A = \begin{pmatrix} 1 & 1 \\ 1 & -1 \end{pmatrix}$, $A\Sigma A^\top = \begin{pmatrix} 11 & 1 \\ 1 & 3 \end{pmatrix}$ (off-diagonal $4 - 3 = 1$).
    \\[0.4em]
    \textbf{Mahalanobis distance} for $\Sigma = \diag(4, 1)$, $\boldsymbol{\mu} = \mathbf{0}$: the point $(2, 1)$ has $d_\Sigma = \sqrt{\frac{4}{4} + \frac{1}{1}} = \sqrt2$, the point $(0, 2)$ has $d_\Sigma = 2$: farther, although closer in Euclidean distance.
}

% ------------------------------------------------------------------ Gaussian
\topic{\secref{sec:gaussian}}{The Gaussian distribution}

\begin{frame}{The multivariate Gaussian}
    \begin{definition}[][def:gaussian]
        \small \(p(\bx) = \dfrac{1}{(2\pi)^{d/2}\det(\Sigma)^{1/2}}\exp\Big(-\frac{1}{2}(\bx - \boldsymbol{\mu})^\top\Sigma^{-1}(\bx - \boldsymbol{\mu})\Big)\), \ $\Sigma$ SPD.
    \end{definition}
    \sidebyside{0.5}{%
        \begin{theorem}[][thm:gaussian-affine]
            \small $\mathbf{Z} \sim \mathcal{N}(\mathbf{0}, I)$, $A$ invertible: \ $A\mathbf{Z} + \boldsymbol{\mu} \sim \mathcal{N}(\boldsymbol{\mu}, AA^\top)$.
        \end{theorem}
        \small
        Sampling: $\boldsymbol{\mu} + L\mathbf{Z}$ with Cholesky $\Sigma = LL^\top$ (Lecture~6). \\[0.3em]
        Level sets: ellipsoids along the eigenvectors of $\Sigma$, semi-axes $\propto\sqrt{\lambda_i}$.
    }{%
        \centering\notesfig[0.72]{L4-gaussian}
    }
    \begin{remark}[The reparameterization trick]
        \small $\boldsymbol{\mu} + \boldsymbol{\sigma} \odot \mathbf{Z}$ is differentiable in $\boldsymbol{\mu}$, $\boldsymbol{\sigma}$: VAEs backpropagate through sampling.
    \end{remark}
\end{frame}

\board{Proof: $A\mathbf{Z} + \boldsymbol{\mu} \sim \mathcal{N}(\boldsymbol{\mu}, AA^\top)$}{Change of variables in $\R^d$: $p_{\mathbf{X}}(\bx) = p_{\mathbf{Z}}(A^{-1}(\bx - \boldsymbol{\mu}))\,|\det A|^{-1}$ (Lecture~7). Then $\norm{A^{-1}(\bx - \boldsymbol{\mu})}^2 = (\bx - \boldsymbol{\mu})^\top\Sigma^{-1}(\bx - \boldsymbol{\mu})$ and $|\det A| = \det(\Sigma)^{1/2}$.}

\begin{frame}{Marginals and conditionals}
    \begin{theorem}[][thm:gaussian-conditional]
        \small $\mathbf{X} = (\mathbf{X}_a, \mathbf{X}_b) \sim \mathcal{N}(\boldsymbol{\mu}, \Sigma)$: \ marginals $\mathbf{X}_a \sim \mathcal{N}(\boldsymbol{\mu}_a, \Sigma_{aa})$, and
        \[\mathbf{X}_a \mid \mathbf{X}_b = \bx_b \ \sim\ \mathcal{N}\big(\boldsymbol{\mu}_a + \Sigma_{ab}\Sigma_{bb}^{-1}(\bx_b - \boldsymbol{\mu}_b),\ \Sigma_{aa} - \Sigma_{ab}\Sigma_{bb}^{-1}\Sigma_{ba}\big).\]
    \end{theorem}
    \centering\fig[0.66]{s09_gaussian_conditional}
    \\[0.1em]
    \raggedright\small
    Mean affine in $\bx_b$; covariance $=$ Schur complement (Lecture~1), $\preceq \Sigma_{aa}$: observing $\mathbf{X}_b$ reduces uncertainty. Galton (1886), \textbf{regression to the mean}: $Y \mid X = x \sim \mathcal{N}(\rho x, 1 - \rho^2)$.
\end{frame}

\board[2]{Proof of the conditional distribution}{Let $K = \Sigma_{ab}\Sigma_{bb}^{-1}$ and $\mathbf{W} = \mathbf{X}_a - K\mathbf{X}_b$. Show $\Cov(\mathbf{W}, \mathbf{X}_b) = 0$, so $\mathbf{W}$ and $\mathbf{X}_b$ are independent (block-diagonal Gaussian), and $\Cov(\mathbf{W}) = \Sigma_{aa} - \Sigma_{ab}\Sigma_{bb}^{-1}\Sigma_{ba}$. Given $\mathbf{X}_b = \bx_b$: $\mathbf{X}_a = \mathbf{W} + K\bx_b$. \\ Example: $\rho = 0.5$, $x = 2$.}

\checkslide{Gaussian conditioning with numbers}{
    \small
    \textbf{Proof steps.} $\Cov(\mathbf{W}, \mathbf{X}_b) = \Sigma_{ab} - K\Sigma_{bb} = \Sigma_{ab} - \Sigma_{ab} = 0$; \ $\Cov(\mathbf{W}) = \Sigma_{aa} - K\Sigma_{ba} - \Sigma_{ab}K^\top + K\Sigma_{bb}K^\top = \Sigma_{aa} - \Sigma_{ab}\Sigma_{bb}^{-1}\Sigma_{ba}$.
    \\[0.4em]
    \textbf{Example 1} (Galton). $(Y, X)$ standard Gaussians with correlation $\rho = 0.5$: \ $\Sigma_{aa} = \Sigma_{bb} = 1$, $\Sigma_{ab} = 0.5$, \ $K = 0.5$.
    \[Y \mid X = 2 \ \sim\ \mathcal{N}(0 + 0.5 \cdot 2,\ 1 - 0.5 \cdot 1 \cdot 0.5) = \mathcal{N}(1, 0.75).\]
    Parents $2$ sd above average: children on average $1$ sd above, with sd $\sqrt{0.75} \approx 0.87$.
    \\[0.4em]
    \textbf{Example 2.} $\boldsymbol{\mu} = (1, 2)^\top$, $\Sigma = \begin{pmatrix} 4 & 2 \\ 2 & 3 \end{pmatrix}$; observe $X_2 = 5$:
    \[X_1 \mid X_2 = 5 \ \sim\ \mathcal{N}\big(1 + \tfrac{2}{3}(5 - 2),\ 4 - \tfrac{2 \cdot 2}{3}\big) = \mathcal{N}\big(3,\ \tfrac{8}{3}\big).\]
    The variance drops from $4$ to $2.67$ whatever value we observe.
    \\[0.4em]
    \textbf{Sums.} $X \sim \mathcal{N}(1, 4)$, $Y \sim \mathcal{N}(2, 9)$ independent: $X + Y \sim \mathcal{N}(3, 13)$, \ $2X - Y \sim \mathcal{N}(0, 4 \cdot 4 + 9) = \mathcal{N}(0, 25)$.
}

\begin{frame}{Gaussians stay Gaussian}
    \begin{corollary}[][cor:gaussian-linear]
        \small $\mathbf{X} \sim \mathcal{N}(\boldsymbol{\mu}, \Sigma)$, $A$ of full row rank: \ $A\mathbf{X} + \mathbf{b} \sim \mathcal{N}(A\boldsymbol{\mu} + \mathbf{b}, A\Sigma A^\top)$. Independent $\mathcal{N}(\mu_1, \sigma_1^2)$ and $\mathcal{N}(\mu_2, \sigma_2^2)$ add to $\mathcal{N}(\mu_1 + \mu_2, \sigma_1^2 + \sigma_2^2)$.
    \end{corollary}
    \begin{theorem}[Independence of blocks][thm:gaussian-blocks]
        \small For Gaussian vectors, uncorrelated blocks ($\Sigma_{ab} = 0$) are independent.
    \end{theorem}
    \small
    \begin{mlbox}
        \small
        \begin{itemize}
            \item Affine maps, marginals, conditionals of Gaussians are Gaussian: everything is linear algebra on $(\boldsymbol{\mu}, \Sigma)$.
            \item Gaussian processes (Lecture~11): Gaussian regression $=$ conditioning a big Gaussian vector on the observed values.
            \item Kalman filters, linear-Gaussian state space models, probabilistic PCA, and the latent space of VAEs.
        \end{itemize}
    \end{mlbox}
\end{frame}

% ------------------------------------------------------------------ Wrap-up
\begin{frame}{Key takeaways}
    \small
    \begin{itemize}
        \item Bayes: $P(B \mid A) = \frac{P(A \mid B)P(B)}{P(A)}$; don't ignore the prior. Sum rule and product rule.
        \item Densities transform with $\frac{1}{|g'|}$; $F^{-1}(U)$ samples from $F$.
        \item $\E$ is linear, $\Var X = \E[X^2] - \E[X]^2$; Markov, Chebyshev, Chernoff, Jensen.
        \item Covariance $=$ inner product, $|\Corr| \leq 1$; uncorrelated $\not\Rightarrow$ independent. $\E[Y \mid X]$ is the best predictor; law of total variance.
        \item $\Var(z) = nv\E[x^2]$: He $v = \frac{2}{n_{\text{in}}}$, Xavier $\frac{2}{n_{\text{in}} + n_{\text{out}}}$.
        \item $\Cov(A\mathbf{X} + \mathbf{b}) = A\Sigma A^\top$; $\Sigma$ is symmetric PSD; top eigenvector $=$ direction of largest variance.
        \item $\mathcal{N}(\boldsymbol{\mu}, \Sigma) = \boldsymbol{\mu} + A\mathbf{Z}$, $AA^\top = \Sigma$; conditionals: mean $\boldsymbol{\mu}_a + \Sigma_{ab}\Sigma_{bb}^{-1}(\bx_b - \boldsymbol{\mu}_b)$, covariance the Schur complement.
    \end{itemize}
\end{frame}

\begin{frame}{Homework and next time}
    \small
    \begin{itemize}
        \item \textbf{Homework:} Lecture 4, Exercises 3, 5, 8, 10, 14 and 16; Programming Exercise P2.
        \item \textbf{Practice:} Practice 4, sections 1--6.
    \end{itemize}
    \vskip0.6em
    \begin{block}{Next: Probability III}
        The law of large numbers, the central limit theorem, Hoeffding's inequality and the delta method; Monte Carlo, importance sampling and the score-function estimator; Markov chains, MCMC and the Bellman equations; entropy, cross-entropy and the KL divergence.
    \end{block}
\end{frame}

\end{document}
